\section{Additional Statistical Tests}
\label{app:tests}

This appendix runs three further hypothesis tests to check the robustness of the empirical results. The first is a Mincer-Zarnowitz efficiency test. The second set consists of residual diagnostics, including the Ljung-Box test and the ARCH LM test. The third is the Harvey-Leybourne-Newbold small-sample correction of the Diebold-Mariano statistic.

% ─────────────────────────────────────────────────────────────────────────────
\subsection{Mincer-Zarnowitz Forecast Efficiency Test}

The Mincer-Zarnowitz (MZ) test \cite{mincer1969economic} checks whether forecasts
are unbiased and informationally efficient.
For each model the auxiliary regression
\begin{equation}
    y_t = \alpha + \beta\,\hat{y}_t + \varepsilon_t
\end{equation}
is estimated by OLS and the joint null hypothesis $H_0{:}\;\alpha = 0,\;\beta = 1$
is assessed with the conventional OLS $F$-statistic. Because the realised rolling
correlation and forecast errors inherit serial dependence from overlapping windows,
the conventional $F$ reference distribution need not hold. The reported values are
therefore diagnostic coefficient summaries, not HAC-robust tests of forecast
efficiency. A confirmatory MZ analysis should use a heteroskedasticity-and-
autocorrelation-consistent Wald test with a lag choice justified by the window length.

Table~\ref{tab:mz_summary} summarises the percentage of the 24
dependency–window experiments in which each model passes the MZ test,
together with its average $\alpha$ and $\beta$ coefficients.

\begin{table}[H]
    \centering
    \caption{Mincer-Zarnowitz forecast efficiency: percentage of 24 experiments
             not rejected by the conventional OLS $F$ diagnostic ($p \geq 0.05$) and average coefficient estimates.}
    \label{tab:mz_summary}
    \begin{tabular}{lrrrr}
        \toprule
        Model & Efficient (\%) & $\bar{\alpha}$ & $\bar{\beta}$ & $\bar{F}$ \\
        \midrule
        AR(1)         & 70.8 & $+0.0006$ & 0.9976 &   2.0 \\
        HAR         & 70.8 & $+0.0006$ & 0.9974 &   2.4 \\
        Ridge       & 45.8 & $-0.0015$ & 0.9992 &   4.1 \\
        RF          & 16.7 & $-0.0019$ & 1.0140 &  26.2 \\
        Naive (last value) &  8.3 & $+0.0097$ & 0.9706 &  16.3 \\
        GBM         &  8.3 & $-0.0046$ & 1.0226 &  35.0 \\
        ElasticNet  &  4.2 & $-0.0056$ & 1.0175 &  20.9 \\
        Ensemble    &  0.0 & $-0.0077$ & 1.0252 &  30.1 \\
        XGBoost    &  0.0 & $-0.0252$ & 1.0979 & 202.2 \\
        DCC-GARCH  &  0.0 & $-0.0881$ & 1.4433 & 237.3 \\
        \bottomrule
    \end{tabular}
\end{table}

AR(1) and HAR are not rejected by the conventional diagnostic at the same rate
(70.8\% of experiments). Ridge is not rejected in 45.8\% of experiments among
the regularised ML models. These proportions are descriptive because the reported
$F$-tests are not HAC-robust. DCC-GARCH is rejected by the conventional diagnostic
in every experiment; its average $\bar{\beta} = 1.44$ means that realised
correlations swing markedly more than its forecasts do. In the regression
$y = \alpha + \beta\hat{y}$, a slope above one signals forecasts that are too
\emph{compressed}: they systematically under-react to the size of the
movements. This is the over-smoothing of DCC-GARCH seen throughout the main
results.

Table~\ref{tab:mz_gspc_w30} provides the full detail for the representative
pair BTC-USD vs S\&P~500 at $w = 30$.

\begin{table}[H]
    \centering
    \caption{Mincer-Zarnowitz test: BTC-USD vs S\&P~500, $w = 30$.
             Conventional OLS $F$ diagnostic for $H_0{:}\;\alpha=0, \beta=1$;
             $p_\text{MZ}$ is not HAC-robust.}
    \label{tab:mz_gspc_w30}
    \begin{tabular}{lrrrrl}
        \toprule
        Model & $\hat{\alpha}$ & $\hat{\beta}$ & $F$ & $p_\text{MZ}$ & Efficient \\
        \midrule
        Ridge       & $+0.0071$ & 0.9903 &   4.76 & 0.009 & No  \\
        AR(1)         & $+0.0071$ & 0.9909 &   5.07 & 0.006 & No  \\
        HAR         & $+0.0074$ & 0.9913 &   5.77 & 0.003 & No  \\
        Naive (last value) & $+0.0114$ & 0.9695 &  16.74 & $<0.001$ & No  \\
        ElasticNet  & $+0.0065$ & 1.0048 &  17.18 & $<0.001$ & No  \\
        Ensemble    & $+0.0086$ & 1.0113 &  39.91 & $<0.001$ & No  \\
        RF          & $+0.0221$ & 0.9951 &  57.04 & $<0.001$ & No  \\
        GBM         & $+0.0226$ & 0.9945 &  62.30 & $<0.001$ & No  \\
        DCC-GARCH  & $-0.0886$ & 1.4616 & 261.03 & $<0.001$ & No  \\
        XGBoost    & $+0.0436$ & 1.0316 & 318.18 & $<0.001$ & No  \\
        \bottomrule
    \end{tabular}
\end{table}

At this particular pair and window all models are rejected by the conventional
OLS diagnostic, but the coefficient gradient is steep. Ridge ($F = 4.76$), AR(1)
($F = 5.07$), and HAR ($F = 5.77$) have coefficients close to $(0, 1)$ under
this descriptive regression. The DCC-GARCH statistic (261.0) and XGBoost statistic
(318.2) are much larger, but neither difference is treated as formal evidence of
misspecification without HAC-robust inference.

% ─────────────────────────────────────────────────────────────────────────────
\subsection{Residual Diagnostics: Ljung-Box and ARCH LM Tests}

Two classical residual tests are applied to the OOS forecast errors
$e_t = y_t - \hat{y}_t$ for every model $\times$ dependency $\times$ window
combination (240 cases in total).

\textbf{Ljung-Box Q-test} (10 lags) \cite{ljung1978measure}: tests $H_0{:}$ no serial autocorrelation.
Rejection indicates that the model has not fully captured the temporal structure
of the target series.

\textbf{ARCH LM test} (5 lags) \cite{engle1982arch}: regresses $e_t^2$ on its own lags and tests
$H_0{:}$ no conditional heteroscedasticity.
Rejection signals remaining volatility clustering in the residuals.

Table~\ref{tab:resid_summary} reports, for each model, the percentage of
experiments in which the residuals pass each test at the 5\% level.

\begin{table}[H]
    \centering
    \caption{Residual diagnostics: percentage of 24 experiments passing the
             Ljung-Box autocorrelation test and the ARCH LM heteroscedasticity
             test at $\alpha = 0.05$.}
    \label{tab:resid_summary}
    \begin{tabular}{lrr}
        \toprule
        Model & No autocorr (\%) & No ARCH (\%) \\
        \midrule
        Naive (last value) & 50.0 & 58.3 \\
        AR(1)         & 41.7 & 62.5 \\
        HAR         & 41.7 & 58.3 \\
        Ridge       & 20.8 & 54.2 \\
        ElasticNet  &  0.0 & 37.5 \\
        Ensemble    &  0.0 & 12.5 \\
        RF          &  0.0 &  0.0 \\
        GBM         &  0.0 &  0.0 \\
        XGBoost    &  0.0 &  0.0 \\
        DCC-GARCH  &  0.0 &  0.0 \\
        \bottomrule
    \end{tabular}
\end{table}

The autocorrelation results show a clear pattern. The Naive (last-value) model and the autoregressive models (AR(1), HAR) capture part of the serial structure and pass the Ljung-Box test in 42--50\% of experiments. Ridge passes in 21\% of cases, showing that it picks up the autoregressive signal from the lag features. ElasticNet is a particular case: it fails every autocorrelation test (0\%) yet still has a meaningful ARCH pass rate (38\%), so it removes some temporal dependence but not all. The tree-based methods (RF, GBM, XGBoost) and DCC-GARCH retain significant autocorrelation and heteroscedasticity in every single experiment: they never encode the target's autoregressive structure, so that structure leaks into the residuals, which remain far from white noise. The large Q-statistics for the tree ensembles (for example, $Q > 1{,}000$ for XGBoost on most pairs) make this clear: these forecasters simply do not carry the AR structure forward from training.

% ─────────────────────────────────────────────────────────────────────────────
\subsection{Harvey-Leybourne-Newbold Corrected Diebold-Mariano Test}

The standard Diebold-Mariano statistic uses the asymptotic standard-normal
distribution, which tends to over-reject in finite samples.
Harvey, Leybourne \& Newbold \cite{harvey1997testing} propose multiplying
the DM statistic by the correction factor
\begin{equation}
    k = \sqrt{\frac{n + 1 - 2h + h(h-1)/n}{n}},
\end{equation}
where $n$ is the number of forecast observations and $h$ the forecast horizon,
and comparing the result to a $t(n-1)$ distribution.
For $h = 1$ the factor simplifies to $k = \sqrt{(n-1)/n}$, which is slightly
below one and therefore shrinks the statistic; with $n > 1{,}000$ the shrinkage
is of the order of $10^{-4}$. This correction addresses finite-sample size, not
the serial dependence in loss differentials induced by overlapping rolling windows.

Applied to all 336 DM comparisons (168 ML models vs DCC-GARCH and 168 vs
Naive (last value)), Table~\ref{tab:dm_hln} shows that the HLN correction changes
the reported zero-lag statistics only trivially (average correction factor
$k = 0.999774$). This establishes only that the displayed calculations are
insensitive to the finite-sample HLN multiplier. It does not validate their
zero-lag variance estimator or turn them into formal significance tests for the
overlapping rolling target.

\begin{table}[H]
    \centering
    \caption{HLN-corrected DM diagnostic: ML models vs DCC-GARCH benchmark,
             BTC-USD vs S\&P~500, $w = 30$.
             DM: original statistic (normal); DM$_\text{HLN}$: corrected
             statistic ($t_{n-1}$).
             Reported $p$-values retain the zero-lag variance assumption and are descriptive.}
    \label{tab:dm_hln}
    \begin{tabular}{lrrrrl}
        \toprule
        Model & DM & $p_\text{DM}$ & DM$_\text{HLN}$ & $p_\text{HLN}$ & Sig. \\
        \midrule
        Ridge      & 26.86 & $<0.001$ & 26.86 & $<0.001$ & *** \\
        HAR        & 26.83 & $<0.001$ & 26.83 & $<0.001$ & *** \\
        ElasticNet & 26.97 & $<0.001$ & 26.97 & $<0.001$ & *** \\
        Ensemble   & 27.05 & $<0.001$ & 27.04 & $<0.001$ & *** \\
        RF         & 26.44 & $<0.001$ & 26.43 & $<0.001$ & *** \\
        GBM        & 26.16 & $<0.001$ & 26.16 & $<0.001$ & *** \\
        XGBoost   & 24.29 & $<0.001$ & 24.28 & $<0.001$ & *** \\
        \bottomrule
    \end{tabular}
\end{table}

Two aspects of the comparison set in Table~\ref{tab:dm_hln} should be made explicit.
First, the seven \emph{fitted} models are compared against two benchmarks: the DCC-GARCH
econometric model and the naive last-value baseline. AR(1) is treated as one of the
persistence baselines rather than as a test model, so it does not appear as a separate
row; its position relative to the fitted models is established instead by the
average-RMSE ranking of Table~\ref{tab:model_ranking_summary}, where AR(1), HAR, and
Ridge have closely clustered errors. Second, the Diebold--Mariano test assumes
non-nested forecasts, whereas every fitted model here includes the lagged target among
its predictors and therefore nests the persistence baselines. The statistics against
Naive (last value), and by extension any test against AR(1), should accordingly be read as
descriptive; the conclusions rest on the comparison the test is designed for, namely the
fitted models against the non-nested DCC-GARCH benchmark.

% ─────────────────────────────────────────────────────────────────────────────
\subsection{Automated Validation Suite}
\label{subsec:auto_tests}

In addition to the statistical tests above, the full project ships with an
automated \texttt{pytest} suite of \textbf{81 unit and integration tests}
organised into ten test classes. These are predominantly consistency and regression
checks: they verify that a fresh run reproduces the published outputs (file presence,
value ranges, sign conventions, and key relationships such as the average ranking of
AR(1) against DCC-GARCH); they do not prove the methods correct in an absolute sense.
The suite is executed automatically as part of the reproducibility runner
(\texttt{python run\_all.py}) between the notebook step and the final
\LaTeX\ compilation, so every rebuilt PDF contains a fresh test report.

The test classes cover the following aspects of the pipeline:
\begin{itemize}
    \item \textbf{Dataset Integrity}. Price and return files exist. They contain the correct tickers, positive prices, sorted and non-duplicated dates, and extend to at least 2026.
    \item \textbf{Data Quality}. Missing-rate thresholds are met, zero-price detection works, annualised volatility is reasonable, and pairwise coverage overlap is adequate.
    \item \textbf{Feature Engineering}. The Fisher-$z$ round-trip identity holds, rolling correlation falls within the expected range, and prediction CSVs contain no look-ahead information.
    \item \textbf{Model Metrics}. RMSE is greater than zero, $R^2$ is less than or equal to one, MAE is less than or equal to RMSE, all windows and expected models are present, and AR(1) performs better than DCC-GARCH on average.
    \item \textbf{DM Tests}. DM statistics are finite, $p$-values lie between zero and one, and the sign convention for Ridge versus DCC-GARCH is correct.
    \item \textbf{Signal Metrics}. Balanced accuracy and F1 score are between zero and one, AUC is above random, and the exit rate falls within a sensible range.
    \item \textbf{Walk-Forward}. The correct number of prediction CSVs are produced, predictions are non-trivial, and indices align with returns.
    \item \textbf{Output Files}. After a full pipeline run, all expected CSV, JSON, PNG, and \texttt{.tex} artefacts are present.
    \item \textbf{Bootstrap CI}. The lower bound is less than or equal to the mean, which is less than or equal to the upper bound for RMSE and $R^2$, and CI width is below 0.10.
    \item \textbf{Sensitivity}. All five \texttt{refit\_every} values are present. RMSE lies within a physically plausible range.
\end{itemize}

Table~\ref{tab:test_report} summarises the suite by test class --- the number of
tests, the pass rate, and the total execution time for each. The complete
per-test log, with parametrised cases and individual timings, is stored in
\texttt{outputs/results/test\_results.xml} in the repository. All 81 tests pass
on the present outputs, with none skipped.

\needspace{4\baselineskip}

\IfFileExists{tables/test_report_auto.tex}{%
    \begin{longtable}{p{3.2cm}p{6.5cm}cr}
\caption{Automated validation suite: 81 tests covering dataset integrity, feature engineering, model metrics, DM tests, signal layer, and reproducibility.}
\label{tab:test_report} \\
\toprule
Test class & Description & Result & s \\
\midrule
\endfirsthead
\multicolumn{4}{l}{\small\textit{(continued from previous page)}} \\
\toprule
Test class & Description & Result & s \\
\midrule
\endhead
\midrule
\multicolumn{4}{r}{\small\textit{continued on next page}} \\
\endfoot
\bottomrule
\endlastfoot
TestDatasetIntegrity & prices file exists & \textcolor{teal}{Pass} & 0.001 \\
 & returns file exists & \textcolor{teal}{Pass} & 0.000 \\
 & prices has all tickers & \textcolor{teal}{Pass} & 0.012 \\
 & prices all positive & \textcolor{teal}{Pass} & 0.008 \\
 & prices no all nan row & \textcolor{teal}{Pass} & 0.006 \\
 & prices index sorted & \textcolor{teal}{Pass} & 0.006 \\
 & prices no duplicate dates & \textcolor{teal}{Pass} & 0.006 \\
 & prices minimum length & \textcolor{teal}{Pass} & 0.006 \\
 & returns matches prices columns & \textcolor{teal}{Pass} & 0.012 \\
 & returns are log returns & \textcolor{teal}{Pass} & 0.007 \\
 & returns length is prices minus one & \textcolor{teal}{Pass} & 0.012 \\
 & eth start date correct & \textcolor{teal}{Pass} & 0.006 \\
 & prices end date reasonable & \textcolor{teal}{Pass} & 0.006 \\
TestDataQuality & missing rate below threshold & \textcolor{teal}{Pass} & 0.007 \\
 & no negative prices & \textcolor{teal}{Pass} & 0.007 \\
 & no zero prices & \textcolor{teal}{Pass} & 0.007 \\
 & return volatility reasonable & \textcolor{teal}{Pass} & 0.007 \\
 & coverage overlap matrix symmetric & \textcolor{teal}{Pass} & 0.002 \\
 & btc eth full overlap & \textcolor{teal}{Pass} & 0.007 \\
 & returns mean near zero & \textcolor{teal}{Pass} & 0.009 \\
 & no constant price series & \textcolor{teal}{Pass} & 0.010 \\
 & data quality stats file & \textcolor{teal}{Pass} & 0.002 \\
TestFeatureEngineering & fisher z identity & \textcolor{teal}{Pass} & 0.001 \\
 & fisher z clips boundaries & \textcolor{teal}{Pass} & 0.002 \\
 & rolling corr range & \textcolor{teal}{Pass} & 0.011 \\
 & rolling corr window effect & \textcolor{teal}{Pass} & 0.010 \\
 & no lookahead in predictions & \textcolor{teal}{Pass} & 0.033 \\
 & target aligned with features & \textcolor{teal}{Pass} & 0.008 \\
TestModelMetrics & metrics file exists & \textcolor{teal}{Pass} & 0.000 \\
 & metrics has expected columns & \textcolor{teal}{Pass} & 0.002 \\
 & rmse positive & \textcolor{teal}{Pass} & 0.002 \\
 & r2 below one & \textcolor{teal}{Pass} & 0.002 \\
 & mae leq rmse & \textcolor{teal}{Pass} & 0.002 \\
 & all windows present & \textcolor{teal}{Pass} & 0.002 \\
 & expected models present & \textcolor{teal}{Pass} & 0.002 \\
 & ar1 beats dcc on average & \textcolor{teal}{Pass} & 0.003 \\
 & naive last close to ar1 & \textcolor{teal}{Pass} & 0.002 \\
 & longer window lower rmse & \textcolor{teal}{Pass} & 0.011 \\
 & n consistent & \textcolor{teal}{Pass} & 0.005 \\
 & high r2 for longer windows & \textcolor{teal}{Pass} & 0.003 \\
TestDMTests & dm file exists & \textcolor{teal}{Pass} & 0.000 \\
 & dm has expected columns & \textcolor{teal}{Pass} & 0.002 \\
 & dm stat finite & \textcolor{teal}{Pass} & 0.002 \\
 & p value in unit interval & \textcolor{teal}{Pass} & 0.002 \\
 & ridge beats dcc dm positive & \textcolor{teal}{Pass} & 0.002 \\
 & ridge vs naive dm mostly negative & \textcolor{teal}{Pass} & 0.002 \\
 & most dcc comparisons significant & \textcolor{teal}{Pass} & 0.003 \\
 & n sufficient for dm & \textcolor{teal}{Pass} & 0.002 \\
TestSignalMetrics & signal file exists & \textcolor{teal}{Pass} & 0.000 \\
 & signal has expected columns & \textcolor{teal}{Pass} & 0.002 \\
 & balanced accuracy range & \textcolor{teal}{Pass} & 0.002 \\
 & f1down range & \textcolor{teal}{Pass} & 0.002 \\
 & auc above random & \textcolor{teal}{Pass} & 0.002 \\
 & exit rate not trivial & \textcolor{teal}{Pass} & 0.003 \\
 & equity signal better than commodities & \textcolor{teal}{Pass} & 0.009 \\
TestWalkForward & prediction csvs exist & \textcolor{teal}{Pass} & 0.000 \\
 & prediction has y true column & \textcolor{teal}{Pass} & 0.020 \\
 & predictions not all nan & \textcolor{teal}{Pass} & 0.022 \\
 & gspc w30 predictions exist & \textcolor{teal}{Pass} & 0.000 \\
 & predictions index matches returns & \textcolor{teal}{Pass} & 0.014 \\
TestOutputFiles & output file exists[D:\\clear\\outputs\\results\\metrics.csv] & \textcolor{teal}{Pass} & 0.000 \\
 & output file exists[D:\\clear\\outputs\\results\\dm tests.csv] & \textcolor{teal}{Pass} & 0.000 \\
 & output file exists[D:\\clear\\outputs\\results\\signal metrics.cs & \textcolor{teal}{Pass} & 0.000 \\
 & output file exists[D:\\clear\\outputs\\results\\run metadata.json & \textcolor{teal}{Pass} & 0.000 \\
 & output file exists[D:\\clear\\outputs\\results\\dataset summary.j & \textcolor{teal}{Pass} & 0.000 \\
 & output file exists[D:\\clear\\outputs\\tables\\metrics table.tex] & \textcolor{teal}{Pass} & 0.000 \\
 & output file exists[D:\\clear\\outputs\\tables\\dm tests.tex] & \textcolor{teal}{Pass} & 0.000 \\
 & output file exists[D:\\clear\\outputs\\tables\\signal metrics.tex & \textcolor{teal}{Pass} & 0.000 \\
 & output file exists[D:\\clear\\outputs\\figures\\dm heatmap Ridge  & \textcolor{teal}{Pass} & 0.001 \\
 & output file exists[D:\\clear\\outputs\\figures\\dataset prices.pn & \textcolor{teal}{Pass} & 0.000 \\
 & output file exists[D:\\clear\\outputs\\figures\\model rmse compar & \textcolor{teal}{Pass} & 0.000 \\
 & metadata json valid & \textcolor{teal}{Pass} & 0.002 \\
 & metrics table tex not empty & \textcolor{teal}{Pass} & 0.000 \\
 & forecast figures all windows & \textcolor{teal}{Pass} & 0.000 \\
TestBootstrapCI & ci file exists & \textcolor{teal}{Pass} & 0.000 \\
 & ci lower leq mean leq upper & \textcolor{teal}{Pass} & 0.002 \\
 & rmse ci width reasonable & \textcolor{teal}{Pass} & 0.002 \\
TestSensitivity & sensitivity file exists & \textcolor{teal}{Pass} & 0.000 \\
 & sensitivity has expected refit values & \textcolor{teal}{Pass} & 0.001 \\
 & sensitivity rmse in range & \textcolor{teal}{Pass} & 0.001 \\
 & sensitivity models present & \textcolor{teal}{Pass} & 0.001 \\
\multicolumn{4}{r}{\small Passed: \textbf{81}\quad Failed: \textbf{0}\quad Skipped: \textbf{0}\quad Total: \textbf{81}} \\
\end{longtable}%
}{%
    \begin{table}[H]
        \centering
        \caption{Automated validation suite results}
        \label{tab:test_report}
        \small Run \texttt{python run\_all.py} to regenerate this table.
    \end{table}%
}

\clearpage
